\documentclass[12pt,a4paper]{article}
\usepackage[margin=25mm, headheight=15pt, headsep=10pt]{geometry}
\usepackage{fontspec}
\usepackage[table]{xcolor} % Note: table option is required for rowcolors
\usepackage{titlesec}
\usepackage{tocloft}
\usepackage{fancyhdr}
\usepackage{booktabs}
\usepackage{tcolorbox}
\tcbuselibrary{skins, breakable}
\usepackage[colorlinks=true, linkcolor=emerald, urlcolor=emerald, bookmarks=true]{hyperref}

\definecolor{emerald}{RGB}{0,102,51}
\definecolor{darkred}{RGB}{139,0,0}
\definecolor{lightgray}{RGB}{245,245,245}

\usepackage[nil]{babel}
\babelprovide[import, main]{bengali}
\babelprovide[import]{arabic}
\babelfont{rm}[Renderer=HarfBuzz,Scale=1.0]{Noto Serif Bengali}
\babelfont{sf}[Renderer=HarfBuzz]{Noto Sans Bengali}
\babelfont{tt}[Renderer=HarfBuzz]{Noto Sans Bengali}
\babelfont[arabic]{rm}[Renderer=HarfBuzz,Scale=1.1]{Noto Naskh Arabic}

% Section styling
\titleformat{\section}{\Large\bfseries\color{emerald}}{\thesection.}{0.5em}{}
\titleformat{\subsection}{\large\bfseries\color{emerald!85!black}}{\thesubsection.}{0.5em}{}
\titleformat{\subsubsection}{\normalsize\bfseries\color{emerald!70!black}}{\thesubsubsection.}{0.5em}{}
\titlespacing*{\section}{0pt}{18pt}{8pt}

% Fancy Header/Footer setup
\pagestyle{fancy}
\fancyhf{} % clear all header and footer fields
\fancyhead[L]{\color{emerald}\small\textbf{\nouppercase{\leftmark}}}
\fancyfoot[L]{\scriptsize\color{black!50}{\lat{AI}}-সংকলিত, আলেম-যাচাইকৃত নয়}
\fancyfoot[C]{\color{emerald!70!black}\thepage}
\renewcommand{\headrulewidth}{0.4pt}
\renewcommand{\headrule}{\hbox to\headwidth{\color{emerald}\leaders\hrule height \headrulewidth\hfill}}
\renewcommand{\footrulewidth}{0pt}

% Custom boxes
% 1. Ayat/Hadith Box
\newtcolorbox{ayatbox}[1][]{
    enhanced, breakable, 
    colback=emerald!5, 
    colframe=emerald, 
    boxrule=0.5pt, 
    arc=3pt, 
    left=10pt, right=10pt, top=10pt, bottom=10pt,
    #1
}

% 2. Quote/Translation Box
\newtcolorbox{quotebox}[1][]{
    enhanced, breakable,
    frame hidden,
    colback=lightgray,
    borderline west={3pt}{0pt}{emerald},
    left=15pt, right=10pt, top=10pt, bottom=10pt,
    sharp corners,
    #1
}

% 3. AI-generated notice box
\newtcolorbox{warnbox}[1][]{
    enhanced, breakable,
    colback=red!4,
    colframe=darkred,
    boxrule=0.8pt,
    arc=3pt,
    left=10pt, right=10pt, top=10pt, bottom=10pt,
    #1
}

% Commands
\newcommand{\arab}[1]{\foreignlanguage{arabic}{#1}}
\newcommand{\kib}[1]{\textcolor{emerald}{\textbf{#1}}}

% Latin (English) fallback: Noto Serif Bengali has NO Latin glyphs
\newfontfamily\latfont[Renderer=HarfBuzz]{Noto Serif}
\newcommand{\lat}[1]{{\latfont #1}}

\setlength{\parskip}{5pt}
\setlength{\parindent}{0pt}

% Fix for missing bullet in Noto Serif Bengali
\renewcommand{\labelitemi}{$\bullet$}



\begin{document}
\begin{center}{\Large\bfseries\color{emerald} পার্ট ৪ — নামাজের সঠিক পদ্ধতি (ধাপে ধাপে, দলিলসহ)}\end{center}
\vspace{1em}
\section*{পার্ট ৪ — নামাজের সঠিক পদ্ধতি (ধাপে ধাপে, দলিলসহ)}

\begin{warnbox}
 \textbf{গুরুত্বপূর্ণ নোটিশ (\lat{Important} \lat{Notice}):} এই কন্টেন্টটি \lat{AI} (কৃত্রিম বুদ্ধিমত্তা) দিয়ে প্রস্তুত ও সংকলিত। \lat{AI} কঠোর সূত্র-মান নীতি মেনে শুধু নির্ভরযোগ্য উৎস থেকে তথ্য সংগ্রহ করেছে — কেবল সহীহ/হাসান হাদিস ও সহীহ সনদের আসার রাখা হয়েছে, দুর্বল সনদ বাদ দেওয়া হয়েছে। তবে এটি কোনো আলেম বা মুহাদ্দিস কর্তৃক যাচাইকৃত নয়।
\end{warnbox}


\begin{quote}
\textbf{পড়ার নিয়ম:} প্রতিটি ধাপে — \textbf{বর্ণনা} $\rightarrow$ \textbf{দলিল} (হাদিস নম্বরসহ) $\rightarrow$ \textbf{মাযহাব-পার্থক্য} (যেখানে আছে) $\rightarrow$ \textbf{সাধারণ ভুল}। যেসব বিষয়ে মাযহাবগুলোর মতভেদ আছে, সেগুলোর বিস্তারিত বিশ্লেষণ পার্ট ৬ (`\lat{09}-বিতর্কিত-মাসআলা-মাযহাব-তুলনা.\lat{md}`) ফাইলে আছে — এখানে সংক্ষেপে ইঙ্গিত দেওয়া হলো।
\end{quote}


\medskip\hrule\medskip

\subsection*{ভূমিকা — পদ্ধতির মূল ভিত্তি}

নামাজের পদ্ধতির সবচেয়ে বড় দলিল তিনটি:

\begin{enumerate}
\item \textbf{\lat{“}তোমরা নামাজ পড়ো, যেমন দেখেছ আমাকে নামাজ পড়তে।\lat{”}} — সহীহ বুখারি ৬৩১ (মালিক ইবনে হুওয়াইরিস)
\item \textbf{মুসি' সালাতুহু হাদিস:} এক ব্যক্তি ঠিকমতো নামাজ পড়তে পারেননি; নবীজি (সা.) তাকে তাকবীর, কিরাআত, রুকু, কাউমা, সিজদা, জালসা — প্রতিটি রুকনে স্থিরতা (ইতমিনান) শিখিয়ে বললেন — \lat{“}তারপর তোমার পুরো নামাজে এরকম করো।\lat{”} — সহীহ বুখারি ৭৫৭; সহীহ মুসলিম ৩৯৭
\item \textbf{আবু হুমাইদ আস-সাঈদী (রাঃ)-এর পূর্ণ বর্ণনা} — দশজন সাহাবির সামনে নবীজির নামাজের প্রতিটি ভঙ্গি — সহীহ বুখারি ৮২৮; সুনানে আবু দাউদ ৭৩০
\end{enumerate}
এ তিনটি দলিলের আলোকে নামাজের ধাপগুলো নিচে বর্ণনা করা হলো। \textbf{নিয়ম:} প্রতিটি ভঙ্গিতে দেহ স্থির রাখা (ইতমিনান) ফরজ — \lat{“}যে রুকু-সিজদায় পিঠ সোজা করে স্থির থাকে না, তার নামাজ বাতিল\lat{”} (সহীহ বুখারি ৭৯১; বুখারি ৭৫৭-এর শিক্ষা)।

\medskip\hrule\medskip

\subsection*{ধাপ ১ — নিয়ত (\lat{Intention})}

\textbf{বর্ণনা:} নামাজের আগে অন্তরে স্থির করা — কোন নামাজ পড়ছি, ফরজ/সুন্নত, কত রাকাত। নিয়তের স্থান \textbf{হৃদয় (\lat{heart})}।

\textbf{দলিল:} \lat{“}নিশ্চয়ই আমলসমূহ নিয়তের উপর নির্ভরশীল...\lat{”} — সহীহ বুখারি ১; সহীহ মুসলিম ১৯০৭।

\textbf{মাযহাব-পার্থক্য:}
\begin{itemize}
\item \textbf{জমহুর (হানাফি, মালিকি, হাম্বলি):} নিয়ত মুখে উচ্চারণ করা সুন্নত নয়; নামাজের নিয়ত মুখে বলার প্রচলন বিদআতের পর্যায়ে — নিয়ত অন্তরের কাজ।
\item \textbf{শাফি'ঈ মাযহাব:} নিয়ত মুখে বলা মুস্তাহাব (তাকবীরের সাথে মিলিয়ে) — তবে এটিও মূলত অন্তরের নিয়তের সহায়ক মাত্র।
\item \textbf{সব মাযহাবের ঐকমত্য:} নিয়তের আসল স্থান অন্তর; মুখে বলা নিয়তের \textbf{বিকল্প (\lat{substitute})} নয়।
\end{itemize}
\textbf{সাধারণ ভুল:} নিয়তের জন্য জোর করে মুখস্থ বাক্য বলা বা ভুলে গেলে নামাজ নষ্ট হওয়ার ভয় — এটি ভুল ধারণা; অন্তরে সংকল্পই যথেষ্ট।

\medskip\hrule\medskip

\subsection*{ধাপ ২ — তাকবীরে তাহরিমা (\lat{Opening} \lat{Takbir})}

\textbf{বর্ণনা:} কিবলামুখী হয়ে দাঁড়িয়ে \lat{“}আল্লাহু আকবার\lat{”} বলে নামাজ শুরু করা। তাকবীরের সাথে হাত কাঁধ/কানের সমান করে ওঠানো (রফউল ইয়াদাইন)।

\textbf{দলিল:}
\begin{itemize}
\item \lat{“}নামাজের চাবি পবিত্রতা, এর তাহরিম (নিষিদ্ধকারী) তাকবীর, আর এর তাহলিল (সমাপ্তিকারী) সালাম।\lat{”} — আবু দাউদ ৬১; তিরমিযী ৪; ইবনে মাজাহ ২৭৫ (সহীহ)
\item \lat{“}নবীজি নামাজ শুরু করার সময়, রুকুতে যাওয়ার সময় ও রুকু থেকে মাথা তোলার সময় হাত উঠাতেন।\lat{”} — বুখারি ৭৩৫-৭৩৭, ৭৩৯; মুসলিম ৩৯০ (সহীহ) — হাতের উচ্চতা: কাঁধ পর্যন্ত (বুখারি ৭৩৫) ও কানের লতি পর্যন্ত (আবু দাউদ ৭২৬, ওয়ায়িলের বর্ণনা) — দুটোই সহীহ
\item \textbf{মাযহাব-পার্থক্য:} তাকবীরে তাহরিমায় হাত ওঠানো সবাই করে; তবে \textbf{রুকুর সময় ও রুকু থেকে ওঠার সময় হাত ওঠানো} নিয়ে মতভেদ — হানাফি ও মালিকি মাযহাবে সুন্নত নয়, শাফি'ঈ ও হাম্বলি মাযহাবে সুন্নত। বিস্তারিত: পার্ট ৬।
\end{itemize}
\textbf{সাধারণ ভুল:} তাকবীরের আগে নড়াচড়া করা, তাকবীর বলার পরও হাত নামাতে দেরি, তাকবীরের সময় দৃষ্টি উপরে তোলা (দৃষ্টি সিজদার স্থানে রাখা সুন্নত)।

\medskip\hrule\medskip

\subsection*{ধাপ ৩ — সানা (ইস্তিফতাহর দোয়া)}

\textbf{বর্ণনা:} তাকবীরের পর হাত বাঁধা অবস্থায় নীরবে সানা পড়া:

\begin{quote}
\lat{“}সুবহানাকা আল্লাহুম্মা ওয়া বিহামদিকা, ওয়া তাবারাকাসমুকা, ওয়া তা'আলা জাদ্দুকা, ওয়া লা ইলাহা গাইরুকা।\lat{”} (হে আল্লাহ! আপনার প্রশংসাসহ আপনি পবিত্র; আপনার নাম বরকতময়, আপনার মহিমা সুউচ্চ; আপনি ছাড়া কোনো উপাস্য নেই।)
\end{quote}

\textbf{দলিল:} আবু দাউদ ৭৭৫; তিরমিযী ২৪৩; ইবনে মাজাহ ৮০৬ (সহীহ, আল-আলবানী)।

\textbf{মাযহাব-পার্থক্য:} সানার পঠন সব মাযহাবেই সুন্নত; শব্দগত ভিন্নতা আছে (যেমন \lat{“}ওয়াজ্জাহতু ওয়াজহিয়া\lat{”} সংস্করণ — মুসলিম ৭৭১-এ আছে), কোনোটি বাধ্যতামূলক নয়। হানাফি মাযহাবে সানা \textbf{ওয়াজিব} পর্যায়ের সুন্নত হিসেবে গণ্য।

\textbf{সাধারণ ভুল:} সানা ভুলে গেলে বা ছেড়ে দিলে নামাজ নষ্ট হয় — এমন ভুল ধারণা; সানা সুন্নত, ছুটলে নামাজ শুদ্ধ, সাহু সিজদার প্রয়োজন নেই (জমহুর)।

\medskip\hrule\medskip

\subsection*{ধাপ ৪ — আউযুবিল্লাহ ও বিসমিল্লাহ}

\textbf{বর্ণনা:} ফাতিহার আগে নীরবে \lat{“}আউযুবিল্লাহি মিনাশ শাইতানির রাজিম\lat{”} ও \lat{“}বিসমিল্লাহির রাহমানির রাহিম\lat{”} পড়া।

\textbf{দলিল:}
\begin{itemize}
\item \lat{“}তুমি যখন কুরআন তিলাওয়াত করো, তখন বিতাড়িত শয়তান থেকে আল্লাহর কাছে আশ্রয় চাও।\lat{”} — সূরা আন-নাহল ১৬:৯৮
\item নবীজি (সা.) বিসমিল্লাহ দিয়ে কিরাআত শুরু করতেন — সহীহ মুসলিম ৫৯৯ (আনাসের বর্ণনা); বুখারি ৭৩৬-এর বর্ণনাতেও ইঙ্গিত আছে
\end{itemize}
\textbf{মাযভাব-পার্থক্য:} বিসমিল্লাহ \textbf{জোরে বলা} নিয়ে মতভেদ — হানাফি, মালিকি, হাম্বলি (প্রসিদ্ধ মত) ও শাফি'ঈ মাযহাবের একটি মত অনুযায়ী জাহরি নামাজে বিসমিল্লাহ আস্তে বলা হয়; শাফি'ঈ মাযহাবের প্রসিদ্ধ মত — জাহরি নামাজে বিসমিল্লাহ জোরে। বিস্তারিত: পার্ট ৬।

\textbf{সাধারণ ভুল:} আউযুবিল্লাহ-বিসমিল্লাহকে ফরজ ভাবা — এগুলো সুন্নত; ভুলে গেলে নামাজ শুদ্ধ।

\medskip\hrule\medskip

\subsection*{ধাপ ৫ — সূরা ফাতিহা (প্রতিটি রাকাতে)}

\textbf{বর্ণনা:} প্রতিটি রাকাতে সূরা ফাতিহা পড়া — নামাজের অন্যতম \textbf{রুকন (\lat{pillar})}; এরপর নীরবে \lat{“}আমিন\lat{”}।

\textbf{দলিল:}
\begin{itemize}
\item \lat{“}যে ব্যক্তি সূরা ফাতিহা পড়েনি, তার নামাজ নেই (পূর্ণ হয়নি)।\lat{”} — বুখারি ৭৫৬; মুসলিম ৩৯৪ (সহীহ)
\item \lat{“}যখন ইমাম ‘গাইরিল মাগদুবি আলাইহিম ওয়ালাদ্দাল্লিন’ বলেন, তখন তোমরা আমিন বলো; কারণ ফেরেশতারাও আমিন বলে... যার আমিন ফেরেশতাদের আমিনের সাথে মিলে যাবে, তার আগের গুনাহ মাফ করে দেওয়া হবে।\lat{”} — বুখারি ৭৮০; মুসলিম ৪১০ (সহীহ)
\end{itemize}
\textbf{মাযহাব-পার্থক্য:}
\begin{itemize}
\item \textbf{ইমামের পেছনে ফাতিহা পড়া} — হানাফি: মুক্তাদি পড়বে না (ইমামের কিরাআতই যথেষ্ট); শাফি'ঈ: মুক্তাদির উপর ফাতিহা ফরজ (সব নামাজে); মালিকি: জাহরি নামাজে পড়বে না, সিররি নামাজে পড়বে; হাম্বলি: জাহরি নামাজে ইমামের কিরাআত শুনবে, সিররি নামাজে পড়বে (মাযহাবের প্রসিদ্ধ মত) — বিস্তারিত: পার্ট ৬।
\item \textbf{আমিন জোরে/আস্তে} — শাফি'ঈ ও হাম্বলি: জাহরি নামাজে ইমাম ও মুক্তাদি জোরে আমিন; হানাফি ও মালিকি: আস্তে। বিস্তারিত: পার্ট ৬।
\end{itemize}
\textbf{সাধারণ ভুল:} ফাতিহা তাড়াতাড়ি পড়া, আয়াতের মধ্যে নিঃশ্বাসের বিরতি না দেওয়া, \lat{“}আমিন\lat{”} টেনে দীর্ঘ করা (জমহুরের মতে সুন্নত নয়)।

\medskip\hrule\medskip

\subsection*{ধাপ ৬ — সূরা মিলানো (কিরাআত)}

\textbf{বর্ণনা:} ফরজ নামাজের প্রথম দুই রাকাতে ফাতিহার পর একটি সূরা বা কুরআনের কিছু আয়াত পড়া; তৃতীয়-চতুর্থ রাকাতে শুধু ফাতিহা।

\textbf{দলিল:} \lat{“}নবীজি (সা.) যুহর ও আসরের প্রথম দুই রাকাতে ফাতিহার সাথে আরেকটি সূরা পড়তেন, আর শেষ দুই রাকাতে শুধু ফাতিহা।\lat{”} — বুখারি ৭৬২; মুসলিম ৪৫১ (আবু কাতাদাহ, সহীহ)

\textbf{মাযহাব-পার্থক্য:}
\begin{itemize}
\item হানাফি: প্রথম দুই রাকাতে সূরা মেলানো \textbf{ওয়াজিব}; জাহরি নামাজে ইমাম জোরে পড়বেন, সিররি নামাজে আস্তে।
\item শাফি'ঈ: সূরা মেলানো সুন্নত; প্রথম দুই রাকাতে ফাতিহার পর তা পড়া মুস্তাহাব।
\item মালিকি ও হাম্বলি: সুন্নত (ফজর ও প্রথম দুই রাকাতে)।
\end{itemize}
\textbf{সাধারণ ভুল:} ইমামের পেছনে জোরে সূরা পড়া (ইমামের কিরাআত শোনা সুন্নত — আ'রাফ ৭:২০৪ \lat{“}যখন কুরআন তিলাওয়াত হয়, তখন তা শোনো\lat{”}); মুক্তাদির জোরে পড়া জমহুরের মতে অনুচিত।

\medskip\hrule\medskip

\subsection*{ধাপ ৭ — রুকু (\lat{Bowing})}

\textbf{বর্ণনা:} \lat{“}আল্লাহু আকবার\lat{”} বলে রুকুতে যাওয়া — হাত দিয়ে হাঁটু ধরতে হবে, পিঠ সোজা ও মাথা পিঠের সমান্তরাল, অন্তত ৩ বার \lat{“}সুবহানা রাব্বিয়াল আজিম\lat{”} (মহান রব পবিত্র)।

\textbf{দলিল:}
\begin{itemize}
\item \lat{“}রুকুতে বলো: সুবহানা রাব্বিয়াল আজিম\lat{”} — আবু দাউদ ৮৭০; তিরমিযী ২৬২ (হাসান সহীহ); মুসলিম ৭৭২ (হুযাইফার বর্ণনায় নবীজির দীর্ঘ রুকু)
\item রুকুতে দোয়া: \lat{“}সুবহানাকা আল্লাহুম্মা রাব্বানা ওয়া বিহামদিকা, আল্লাহুম্মাগফিরলি\lat{”} — বুখারি ৭৯৪ (সহীহ); মুসলিম ৪৮৪
\item রুকুর ভঙ্গি — আবু হুমাইদের বর্ণনায়: \lat{“}নবীজি রুকুতে গিয়ে হাত হাঁটুর উপর রাখতেন, পিঠ সোজা রাখতেন\lat{”} — বুখারি ৮২৮; আবু দাউদ ৭৩০
\end{itemize}
\textbf{মাযহাব-পার্থক্য:} রুকুতে যাওয়ার সময় হাত ওঠানো — হানাফি/মালিকি: হাত ওঠানো হয় না; শাফি'ঈ/হাম্বলি: হাত ওঠানো হয়। (পার্ট ৬) রুকুতে তাসবিহের সংখ্যা — কমপক্ষে একবার, উত্তম ৩ বার; কোনো নির্দিষ্ট সংখ্যা ফরজ নয়।

\textbf{সাধারণ ভুল:} রুকুতে পিঠ সোজা না করা, মাথা অতিরিক্ত নিচু/উঁচু করা, তাড়াতাড়ি উঠে যাওয়া (ইতমিনান ভঙ্গ) — \lat{“}যে রুকু-সিজদায় স্থির হয় না, তার নামাজ বাতিল\lat{”} (বুখারি ৭৯১)।

\medskip\hrule\medskip

\subsection*{ধাপ ৮ — কাউমা (রুকু থেকে দাঁড়ানো)}

\textbf{বর্ণনা:} \lat{“}সামিআল্লাহু লিমান হামিদাহ\lat{”} বলে রুকু থেকে সোজা দাঁড়ানো (ইমাম ও একাকী); মুক্তাদি ও একাকী \lat{“}রাব্বানা লাকাল হামদ\lat{”} বলা; দাঁড়িয়ে শরীর পুরোপুরি স্থির করা।

\textbf{দলিল:}
\begin{itemize}
\item \lat{“}নবীজি রুকু থেকে মাথা তুলে বলতেন: সামিআল্লাহু লিমান হামিদাহ, রাব্বানা লাকাল হামদ\lat{”} — বুখারি ৭৯৫ (সহীহ)
\item \lat{“}যখন ইমাম সামিআল্লাহু লিমান হামিদাহ বলেন, তোমরা রাব্বানা লাকাল হামদ বলো\lat{”} — বুখারি ৭৮৯ (সহীহ)
\item \lat{“}নবীজি রুকু থেকে উঠে এমনভাবে দাঁড়াতেন যে, প্রতিটি কশেরুকা নিজের জায়গায় ফিরে আসত\lat{”} — বুখারি ৮২৮; মুসলিম ৪৯৮ (বারা ইবনে আযিব, সহীহ)
\end{itemize}
\textbf{মাযহাব-পার্থক্য:} কাউমা — হানাফি মাযহাবে ওয়াজিব, জমহুরের মতে রুকন; সবার কাছেই এটি নামাজের অপরিহার্য অংশ। মুক্তাদি কি \lat{“}সামিআল্লাহু...\lat{”} বলবে? জমহুরের মতে মুক্তাদি শুধু \lat{“}রাব্বানা লাকাল হামদ\lat{”} — ইমাম ও একাকী উভয়টি বলে।

\textbf{সাধারণ ভুল:} কাউমায় শরীর স্থির না করে সোজা সিজদায় চলে যাওয়া — এটি ইতমিনান ভঙ্গ; নবীজি বলেছেন, রুকু থেকে পূর্ণভাবে দাঁড়ানো ছাড়া নামাজ শুদ্ধ নয় (বুখারি ৭৫৭-এর শিক্ষা)।

\medskip\hrule\medskip

\subsection*{ধাপ ৯ — সিজদা (\lat{Prostration})}

\textbf{বর্ণনা:} \lat{“}আল্লাহু আকবার\lat{”} বলে সিজদায় যাওয়া — সাতটি অঙ্গের উপর: কপাল (নাকসহ), দুই হাতের তালু, দুই হাঁটু, দুই পায়ের আঙুলের ডগা; অন্তত ৩ বার \lat{“}সুবহানা রাব্বিয়াল আ'লা\lat{”} (সর্বোচ্চ রব পবিত্র); সিজদায় দোয়া করা মুস্তাহাব।

\textbf{দলিল:}
\begin{itemize}
\item \lat{“}আমাকে আদেশ করা হয়েছে সাত অঙ্গের উপর সিজদা করতে: কপাল (নাকসহ), দুই হাত, দুই হাঁটু ও দুই পায়ের আঙুল...\lat{”} — বুখারি ৮১২; মুসলিম ৪৯০-৪৯১ (সহীহ)
\item \lat{“}রুকুতে সুবহানা রাব্বিয়াল আজিম, সিজদায় সুবহানা রাব্বিয়াল আ'লা\lat{”} — আবু দাউদ ৮৭০; তিরমিযী ২৬২
\item \lat{“}বান্দা তার রবের সবচেয়ে নিকটবর্তী হয় সিজদার অবস্থায় — তাই সিজদায় বেশি দোয়া করো\lat{”} — মুসলিম ৪৮২ (সহীহ)
\item সিজদায় দোয়া: \lat{“}সুবহানাকা আল্লাহুম্মা রাব্বানা ওয়া বিহামদিকা, আল্লাহুম্মাগফিরলি\lat{”} — বুখারি ৭৯৪
\end{itemize}
\textbf{মাযহাব-পার্থক্য — সিজদায় যাওয়ার পদ্ধতি (হাত আগে না হাঁটু আগে):}
\begin{itemize}
\item \textbf{শাফি'ঈ ও হাম্বলি:} সিজদায় যাওয়ার সময় \textbf{হাত আগে} রাখা, উঠার সময় হাত পরে; দলিল — \lat{“}তোমাদের কেউ যেন উটের মতো না বসে; হাত দুটি হাঁটুর আগে রাখুক\lat{”} (আবু দাউদ ৮৪০; তিরমিযী ২৬৯ — আল-আলবানীর মতে হাসান)।
\item \textbf{হানাফি ও মালিকি:} \textbf{হাঁটু আগে} রাখা; দলিল — ওয়ায়িল ইবনে হুজরের বর্ণনা \lat{“}নবীজি সিজদায় গিয়ে হাঁটু দুটি হাতের আগে রাখতেন\lat{”} (আবু দাউদ ৮৩৯) — তবে এই সনদে দুর্বলতা আছে; এবং \lat{“}উটের মতো না বসা\lat{”}-র ব্যাখ্যা — একসাথে ঝাঁপিয়ে না পড়া।
\item \textbf{বিচার:} হাত-আগে-বর্ণনার দলিল শক্তিশালী (হাসান); হাঁটু-আগে-বর্ণনার সনদ দুর্বল (আল-আলবানী)। তবে বিষয়টি বৈধ ইখতিলাফ — কোনো পক্ষকে ভুল বলা যায় না; নিরাপদ পদ্ধতি — স্থিরভাবে, কোনো এক অঙ্গ প্রথমে রেখে ধীরে সিজদায় যাওয়া। বিস্তারিত: পার্ট ৬।
\end{itemize}
\textbf{সাধারণ ভুল:} সিজদায় পেট উরুর সাথে লেগে থাকা (নবীজি পেট আলাদা রাখতেন — বুখারি ৮২৮), হাত মাটিতে বিছিয়ে দেওয়া (নবীজি হাত কাঁধের সমান্তরাল রাখতেন — বুখারি ৮২৮), পায়ের আঙুল কিবলামুখী না রাখা, সিজদায় কনুই মাটিতে লেগে থাকা।

\medskip\hrule\medskip

\subsection*{ধাপ ১০ — জালসা (দুই সিজদার মাঝের বসা)}

\textbf{বর্ণনা:} \lat{“}আল্লাহু আকবার\lat{”} বলে সিজদা থেকে মাথা তুলে \textbf{ইফতিরাশ} পদ্ধতিতে বসা — বাম পায়ের উপর বসা, ডান পা খাড়া, হাত উরু/হাঁটুর উপর; বসে দোয়া: \lat{“}রাব্বিগফিরলি, রাব্বিগফিরলি\lat{”} (হে আমার রব, আমাকে ক্ষমা করো)।

\textbf{দলিল:}
\begin{itemize}
\item \lat{“}নবীজি দুই সিজদার মাঝে বাম পায়ের উপর বসতেন, ডান পা খাড়া রাখতেন\lat{”} — মুসলিম ৫৭৯-৫৮০ (সহীহ)
\item \lat{“}নবীজি দুই সিজদার মাঝে বলতেন: আল্লাহুম্মাগফিরলি, ওয়ারহামনি, ওয়া'আফিনি, ওয়াহদিনি, ওয়ারজুকনি\lat{”} (হে আল্লাহ! আমাকে ক্ষমা করো, দয়া করো, সুস্থতা দাও, হিদায়াত দাও ও রিজিক দাও) — আবু দাউদ ৮৫০; তিরমিযী ২৮৪; ইবনে মাজাহ ৮৯৭ (হাসান সহীহ, আল-আলবানী)
\end{itemize}
\textbf{মাযহাব-পার্থক্য:} জালসার বসার পদ্ধতি প্রায় সব মাযহাবে এক — ইফতিরাশ; শেষ বৈঠকে (তাশাহহুদে) \textbf{তাওয়াররুক} (বাম পা ডান পায়ের নিচ থেকে বের করে নিতম্বের উপর বসা) — শাফি'ঈ ও হাম্বলি মাযহাবের সুন্নত, হানাফি ও মালিকি ইফতিরাশেই বসেন। উভয় পদ্ধতিই নবীজির আমল থেকে প্রমাণিত (বুখারি ৮২৮ — আবু হুমাইদ উভয় বর্ণনা করেছেন)।

\textbf{সাধারণ ভুল:} জালসা এড়িয়ে দ্রুত দ্বিতীয় সিজদায় চলে যাওয়া — জালসা ও ইতমিনান ফরজের অংশ (বুখারি ৭৫৭)।

\medskip\hrule\medskip

\subsection*{ধাপ ১১ — দ্বিতীয় সিজদা ও দ্বিতীয় রাকাত}

\textbf{বর্ণনা:} দ্বিতীয় সিজদা প্রথমটির মতো — \lat{“}আল্লাহু আকবার\lat{”} বলে সিজদা, তাসবিহ; তারপর \lat{“}আল্লাহু আকবার\lat{”} বলে দাঁড়ানো (কোনো তাকবীর ছাড়া হাত মাটিতে ভর না দিয়ে); দ্বিতীয় রাকাত প্রথম রাকাতের মতোই — ফাতিহা, সূরা, রুকু, সিজদা।

\textbf{দলিল:}
\begin{itemize}
\item \lat{“}নবীজি প্রতিটি উঠা-বসায় তাকবীর বলতেন\lat{”} — বুখারি ৭৮৬ (সহীহ)
\item দ্বিতীয় সিজদা থেকে ওঠার পদ্ধতি — বুখারি ৮২৮ (আবু হুমাইদ): \lat{“}সিজদা থেকে উঠে দাঁড়াতেন\lat{”}; মুসলিম ৫৭৬-এর বর্ণনায় হাত মাটিতে না রেখে ওঠার ইঙ্গিত আছে
\end{itemize}
\textbf{সাধারণ ভুল:} দাঁড়ানোর সময় মাটিতে হাত ভর দেওয়া (জরুরি প্রয়োজনে ছাড়া সুন্নত নয়); দ্বিতীয় রাকাতে নতুন করে সানা পড়া (সানা শুধু প্রথম রাকাতে)।

\medskip\hrule\medskip

\subsection*{ধাপ ১২ — তাশাহহুদ (প্রথম ও শেষ বৈঠক)}

\textbf{বর্ণনা:} দুই রাকাত পর প্রথম তাশাহহুদ; শেষ রাকাতে পূর্ণ তাশাহহুদ। পঠন (ইবনে মাসউদের শিক্ষা):

\begin{quote}
\lat{“}আত্তাহিয়্যাতু লিল্লাহি ওয়াস সালাওয়াতু ওয়াত তাইয়্যিবাত; আসসালামু আলাইকা আইয়ুহান নাবিয়্যু ওয়া রাহমাতুল্লাহি ওয়া বারাকাতুহ; আসসালামু আলাইনা ওয়া আলা ইবাদিল্লাহিস সালিহিন; আশহাদু আল্লা ইলাহা ইল্লাল্লাহু ওয়া আশহাদু আন্না মুহাম্মাদান আবদুহু ওয়া রাসূলুহ।\lat{”}
\end{quote}

\textbf{দলিল:}
\begin{itemize}
\item \lat{“}নবীজি (সা.) আমাকে তাশাহহুদ শিখিয়েছেন, যেমন কুরআনের সূরা শেখাতেন...\lat{”} — বুখারি ৬২৬৫, ৮৩১; মুসলিম ৪০২ (সহীহ)
\item \textbf{আঙুলের ইশারা:} \lat{“}নবীজি তাশাহহুদের বৈঠকে ডান হাত ডান উরুর উপর, বাম হাত বাম উরুর উপর রাখতেন এবং তর্জনী দিয়ে ইশারা করতেন\lat{”} — মুসলিম ৫৭৯-৫৮০ (সহীহ); ওয়ায়িল ইবনে হুজরের বর্ণনায় আঙুল নাড়ানো ও বৃত্ত বানানো — আবু দাউদ ৯৮৯; নাসাঈ ১২৭০ (সহীহ)
\item \textbf{মাযহাব-পার্থক্য:} আঙুল নাড়ানো — হাম্বলি মাযহাব ও অনেক মুহাদ্দিস (আল-আলবানীসহ): নাড়ানো সুন্নত; হানাফি মাযহাব: শুধু উত্তোলন, নাড়ানো নয়; শাফি'ঈ: \lat{“}লা ইলাহা\lat{”}-র সময় উত্তোলন, \lat{“}ইল্লাল্লাহ\lat{”}-এ নামানো; মালিকি: ডান-বাম না করে সামনে ইশারা। সব পক্ষই বৈধ আমল — বিস্তারিত: পার্ট ৬।
\end{itemize}
\textbf{সাধারণ ভুল:} প্রথম তাশাহহুদে দরুদ-দোয়া পড়ে ফেলা (প্রথম তাশাহহুদে শুধু তাশাহহুদ; দরুদ ও দোয়া শেষ বৈঠকে — হানাফি মাযহাবের প্রসিদ্ধ মত), তাশাহহুদের সময় দৃষ্টি এদিক-ওদিক করা।

\medskip\hrule\medskip

\subsection*{ধাপ ১৩ — দরুদ (\lat{Salawat} \lat{on} \lat{the} \lat{Prophet} \arab{ﷺ})}

\textbf{বর্ণনা:} শেষ বৈঠকে তাশাহহুদের পর দরুদ পড়া (দরুদে ইবরাহিম):

\begin{quote}
\lat{“}আল্লাহুম্মা সাল্লি আলা মুহাম্মাদিন ওয়া আলা আলি মুহাম্মাদ, কামা সাল্লাইতা আলা আলি ইবরাহিম; ইন্নাকা হামিদুম মাজিদ। আল্লাহুম্মা বারিক আলা মুহাম্মাদিন ওয়া আলা আলি মুহাম্মাদ, কামা বারাকতা আলা আলি ইবরাহিম; ইন্নাকা হামিদুম মাজিদ।\lat{”}
\end{quote}

\textbf{দলিল:} আবু মাসউদ আল-আনসারী (রাঃ) বলেন — নবীজি (সা.) আমাদের এভাবে দরুদ শিখিয়েছেন — সহীহ বুখারি ৩৩৭০; সহীহ মুসলিম ৪০৬ (সহীহ)।

\textbf{মাযহাব-পার্থক্য:} দরুদ — শাফি'ঈ মাযহাবে ফরজ (শেষ বৈঠকে), অন্যান্য মাযহাবে ওয়াজিব/সুন্নত; সবার কাছেই উত্তম ও গুরুত্বপূর্ণ।

\textbf{সাধারণ ভুল:} দরুদ এড়িয়ে যাওয়া; দরুদের পর দোয়ার আগে সালাম দেওয়া।

\medskip\hrule\medskip

\subsection*{ধাপ ১৪ — শেষ দোয়া ও সালাম}

\textbf{বর্ণনা:} দরুদের পর আল্লাহর কাছে দোয়া (যেমন — \lat{“}আল্লাহুম্মা ইন্নি আউযুবিকা মিন আজাবিল কবরি...\lat{”} — বুখারি ১৩৭৭; মুসলিম ৫৮৮); তারপর ডানে \lat{“}আসসালামু আলাইকুম ওয়া রাহমাতুল্লাহ\lat{”}, বামে \lat{“}আসসালামু আলাইকুম ওয়া রাহমাতুল্লাহ\lat{”} বলে নামাজ শেষ।

\textbf{দলিল:}
\begin{itemize}
\item \lat{“}নবীজি (সা.) সালাম ফিরাতেন ডানে ও বামে — যতক্ষণ না গালের শুভ্রতা দেখা যেত\lat{”} — মুসলিম ৫৮২; বুখারি ৮২৯ (সহীহ)
\item সালামই নামাজের সমাপ্তি — আবু দাউদ ৬১; তিরমিযী ৪
\end{itemize}
\textbf{মাযহাব-পার্থক্য:} সালাম — জমহুরের মতে রুকন (একটি সালামই যথেষ্ট, দুইটি সুন্নত); হানাফি মাযহাবে সালাম ওয়াজিব। শাফি'ঈ মাযহাবে প্রথম সালাম ফরজ, দ্বিতীয়টি সুন্নত।

\textbf{সাধারণ ভুল:} সালামের সময় মাথা না ঘুরিয়ে শুধু চোখ দিয়ে ইশারা করা; সালামের আগে-পরে নফসি কথা বলা।

\medskip\hrule\medskip

\subsection*{ধাপ ১৫ — নামাজের পরের যিকর}

\textbf{বর্ণনা:} সালামের পর — \lat{“}আস্তাগফিরুল্লাহ\lat{”} (৩ বার), \lat{“}আল্লাহুম্মা আনতাস সালামু ওয়া মিনকাস সালাম...\lat{”} (মুসলিম ৫৯১); তারপর ৩৩ বার সুবহানাল্লাহ, ৩৩ বার আলহামদুলিল্লাহ, ৩৩ বার আল্লাহু আকবার এবং ১ বার \lat{“}লা ইলাহা ইল্লাল্লাহু ওয়াহদাহু লা শারিকা লাহ...\lat{”} (মুসলিম ৫৯৬)।

\textbf{দলিল:}
\begin{itemize}
\item সালামের পরের দোয়া — মুসলিম ৫৯১ (সহীহ)
\item ৩৩-৩৩-৩৩ তাসবিহ — মুসলিম ৫৯৬ (সহীহ)
\end{itemize}
\textbf{মাযহাব-পার্থক্য:} যিকরের ধরন-পরিমাণে ভিন্নতা আছে (যেমন হানাফি মাযহাবে যিকর আঙুলে/তাসবিহে গণনা); মূল যিকরগুলো সব মাযহাবেই মুস্তাহাব।

\medskip\hrule\medskip

\subsection*{নারীদের নামাজের পদ্ধতি — সংক্ষিপ্ত আলোচনা}

\textbf{মূল নীতি:} নারীরাও পুরুষের মতোই নামাজের \textbf{মৌলিক (\lat{fundamental})} কাঠামো পালন করে — ফাতিহা, রুকু, সিজদা, তাশাহহুদ — কোনো ভিন্নতা নেই। দলিল — \lat{“}তোমরা নামাজ পড়ো, যেমন দেখেছ আমাকে নামাজ পড়তে\lat{”} (বুখারি ৬৩১) সবার জন্য প্রযোজ্য; কুরআন-সুন্নাহয় নারীদের জন্য আলাদা পদ্ধতির কোনো সহীহ নির্দেশ নেই।

\textbf{যে পার্থক্যগুলো প্রচলিত:}
\begin{itemize}
\item \textbf{সতর:} নারীদের নামাজে মুখ ও হাত ছাড়া পুরো শরীর ঢাকা শর্ত — \lat{“}বালেগা নারীর নামাজ খিমার (মাথা-ঢাকা কাপড়) ছাড়া কবুল হয় না\lat{”} — আবু দাউদ ৬৪১; তিরমিযী ৩৭৭; ইবনে মাজাহ ৬৫৫ (সহীহ, আল-আলবানী)।
\item \textbf{আওয়াজ:} নারীদের নামাজে কণ্ঠ নিচু রাখা; তবে জরুরি প্রয়োজনে (যেমন ইমামের ভুল) \textbf{তাসবিহ নয়, হাততালি} — \lat{“}তাসবিহ পুরুষদের জন্য, হাততালি নারীদের জন্য\lat{”} — বুখারি ১২০৩; মুসলিম ৪২২ (সহীহ)।
\item \textbf{হানাফি মাযহাবের অতিরিক্ত নিয়ম:} সিজদায় শরীর গুটিয়ে রাখা, রুকুতে হাত হাঁটুর উপর রাখা — এগুলো এসেছে কয়েকটি দুর্বল বর্ণনা থেকে; জমহুর (মালিকি, শাফি'ঈ, হাম্বলি) ও ইমাম আহমাদের বক্তব্য — \textbf{নারী-পুরুষের নামাজের ভঙ্গি একই}। বিস্তারিত: পার্ট ৬।
\item \textbf{জামাত:} নারীদের জন্য ঘরে নামাজ উত্তম — \lat{“}নারীদের জন্য সর্বোত্তম নামাজ ঘরে\lat{”} — আবু দাউদ ৫৭০ (সহীহ, আল-আলবানী); তবে মসজিদে জামাতে অংশ নেওয়া জায়েজ — নবীজির যুগে নারীরা মসজিদে আসতেন (বুখারি ৮৬৫)।
\end{itemize}
\medskip\hrule\medskip

\subsection*{পূর্ণ পদ্ধতির চেকলিস্ট (\lat{Checklist})}

\begin{table}[h]\centering\small
\begin{tabular}{lll}
ধাপ & কাজ & ফরজ/সুন্নত \\
\hline
১ & নিয়ত (অন্তরে) & ফরজ (শর্ত) \\
২ & কিবলামুখী দাঁড়ানো & ফরজ (শর্ত) \\
৩ & তাকবীরে তাহরিমা & ফরজ (রুকন) \\
৪ & সানা & সুন্নত \\
৫ & আউযুবিল্লাহ-বিসমিল্লাহ & সুন্নত \\
৬ & সূরা ফাতিহা & ফরজ (রুকন) \\
৭ & সূরা মেলানো (১ম-২য় রাকাত) & ওয়াজিব (হানাফি)/সুন্নত (জমহুর) \\
৮ & রুকু + ইতমিনান + তাসবিহ & ফরজ + সুন্নত \\
৯ & কাউমা (সোজা দাঁড়ানো) & ফরজ (রুকন) \\
১০ & সিজদা (৭ অঙ্গ) + ইতমিনান & ফরজ (রুকন) \\
১১ & জালসা + ইতমিনান & ফরজ (রুকন) \\
১২ & দ্বিতীয় সিজদা & ফরজ (রুকন) \\
১৩ & প্রথম তাশাহহুদ (২ রাকাত পর) & ওয়াজিব (হানাফি)/সুন্নতে মুয়াক্কাদা (জমহুর) \\
১৪ & শেষ বৈঠক: তাশাহহুদ + দরুদ + দোয়া & ফরজ/ওয়াজিব \\
১৫ & সালাম (ডানে-বামে) & ফরজ (রুকন) \\
১৬ & পরের যিকর & মুস্তাহাব \\
\hline
\end{tabular}
\end{table}

\begin{quote}
\textbf{মনে রাখবেন:} ফরজ/ওয়াজিব/সুন্নতের এই শ্রেণিবিভাগে মাযহাবগুলোর নিজস্ব টার্মিনোলজি আছে — বিস্তারিত দেখুন `\lat{03}-ফরজ-ওয়াজিব-সুন্নত-মাকরূহ-শ্রেণিবিভাগ.\lat{md}`।
\end{quote}

\medskip\hrule\medskip
\textit{সংশ্লিষ্ট ফাইল: `হাদিস/\lat{02}-হাদিস-ও-আসার.\lat{md}`, `\lat{09}-বিতর্কিত-মাসআলা-মাযহাব-তুলনা.\lat{md}`}


\end{document}
